\documentclass[a4paper,12pt]{article}
\usepackage[utf8]{inputenc}
\usepackage[spanish]{babel}
\usepackage{listings}
\usepackage{xcolor}
\usepackage{geometry}

\geometry{margin=2.5cm}

\definecolor{codegray}{rgb}{0.5,0.5,0.5}
\definecolor{codepurple}{rgb}{0.58,0,0.82}
\definecolor{backcolour}{rgb}{0.95,0.95,0.92}

\lstset{
    backgroundcolor=\color{backcolour},
    commentstyle=\color{codegray},
    keywordstyle=\color{blue},
    numberstyle=\tiny\color{codegray},
    stringstyle=\color{codepurple},
    basicstyle=\ttfamily\small,
    breakatwhitespace=false,         
    breaklines=true,                 
    captionpos=b,                    
    keepspaces=true,                 
    numbers=left,                    
    numbersep=5pt,                  
    showspaces=false,                
    showstringspaces=false,
    showtabs=false,                  
    tabsize=2
}

\title{Manual de Administración de Sistemas: Terminal y Shell Scripting}
\author{Generado por Inteligencia Artificial (Gemini)}
\date{Febrero 2026}

\begin{document}

\maketitle
\newpage
\tableofcontents
\newpage

\section{Comandos Básicos}
Los comandos básicos permiten la navegación y gestión esencial de archivos en el sistema operativo.

\subsection{Ejemplos de Uso}
\begin{enumerate}
    \item \textbf{Listado de archivos:} El comando \texttt{ls} permite ver el contenido del directorio actual.
    \begin{lstlisting}[language=bash]
    $ ls -lh
    \end{lstlisting}
    \textit{Muestra archivos en formato de lista legible para humanos.}

    \item \textbf{Creación de directorios:} El comando \texttt{mkdir}.
    \begin{lstlisting}[language=bash]
    $ mkdir backup_proyecto
    \end{lstlisting}

    \item \textbf{Copiar archivos:} El comando \texttt{cp}.
    \begin{lstlisting}[language=bash]
    $ cp configuracion.conf configuracion.conf.bak
    \end{lstlisting}
\end{enumerate}

\section{Comandos Avanzados}
Estos comandos son fundamentales para la administración profunda y el monitoreo del sistema.

\subsection{Ejemplos de Uso}
\begin{enumerate}
    \item \textbf{Búsqueda de archivos:} El comando \texttt{find}.
    \begin{lstlisting}[language=bash]
    $ find /var/log -name "*.log" -mtime -7
    \end{lstlisting}
    \textit{Busca archivos .log modificados en los últimos 7 días.}

    \item \textbf{Monitoreo de procesos:} El comando \texttt{top} o \texttt{htop}.
    \begin{lstlisting}[language=bash]
    $ top -u root
    \end{lstlisting}

    \item \textbf{Cambio de permisos:} El comando \texttt{chmod}.
    \begin{lstlisting}[language=bash]
    $ chmod 755 script_servidor.sh
    \end{lstlisting}
\end{enumerate}

\section{Programación Shell (Shell Scripting)}
El Shell Scripting es la automatización de tareas mediante la ejecución de una serie de comandos guardados en un archivo.

\subsection{Estructuras de Control}

\subsubsection{Condicional (If-Else)}
Utilizado para verificar estados del sistema, como si un servicio está activo.
\begin{lstlisting}[language=bash]
#!/bin/bash
SERVICIO="ssh"
if systemctl is-active --quiet $SERVICIO; then
    echo "El servicio $SERVICIO esta funcionando."
else
    echo "Alerta: $SERVICIO esta detenido."
fi
\end{lstlisting}

\subsubsection{Bucle For}
Ideal para realizar tareas repetitivas sobre un conjunto de datos.
\begin{lstlisting}[language=bash]
#!/bin/bash
for usuario in user1 user2 user3; do
    echo "Creando carpeta personal para: $usuario"
    mkdir -p /home/backups/$usuario
done
\end{lstlisting}

\subsubsection{Bucle While}
Útil para monitorear condiciones que cambian con el tiempo.
\begin{lstlisting}[language=bash]
#!/bin/bash
# Monitorear espacio en disco cada 10 segundos
while true; do
    df -h | grep "/dev/sda1"
    sleep 10
done
\end{lstlisting}

\section{Tuberías (Pipes)}
Las tuberías permiten redirigir la salida de un comando para que sirva como entrada de otro, conectando herramientas pequeñas para realizar tareas complejas.

\subsection{Ejemplos}
\begin{itemize}
    \item \textbf{Contar líneas de un log:}
    \begin{lstlisting}[language=bash]
    cat /var/log/auth.log | wc -l
    \end{lstlisting}
    
    \item \textbf{Filtrar y ordenar procesos por memoria:}
    \begin{lstlisting}[language=bash]
    ps aux | sort -nk 4 | tail -n 5
    \end{lstlisting}
    \textit{Muestra los 5 procesos que consumen más memoria RAM.}
\end{itemize}

\section{Errores Comunes y Soluciones}
A continuación, una lista de errores frecuentes en la terminal.

\begin{table}[h!]
\centering
\begin{tabular}{|p{3cm}|p{4cm}|p{5cm}|}
\hline
\textbf{Error} & \textbf{Causa} & \textbf{Solución} \\ \hline
\texttt{Permission denied} & El usuario actual no tiene privilegios suficientes. & Anteponer \texttt{sudo} al comando o revisar permisos con \texttt{ls -l}. \\ \hline
\texttt{Command not found} & El comando no está instalado o no está en el PATH. & Instalar el paquete o verificar la ruta absoluta del binario. \\ \hline
\texttt{No space left on device} & El disco o la partición están llenos. & Eliminar archivos temporales o logs antiguos en \texttt{/var/log}. \\ \hline
\end{tabular}
\end{table}

\section{Transparencia en el uso de IA}
\subsection{Prompt Utilizado}
El presente documento fue generado siguiendo las instrucciones del siguiente prompt: 
\begin{quote}
    \textit{"Generar un manual en base a las instrucciones del archivo adjunto, que contemple lo siguiente: una seccion de comandos basicos incluidos 2 a 3 ejemplos desde la terminal, luego comandos avanzados, tb con 2 a 3 ejemplos. Otra seccion conexplicacion de programacion shell, con ejemplos para cada estructura de programacion, con ejemplos sencillos pero relacionados la administracion del sistema, otra seccion con una lista de errores que nos puede enviar el sistema y cual su solucion y porque ocurre, por ultimo una seccion de tuberias con ejemplos. Siendo totalmente transpatente en el uso de IA dando a conocer en el documento como otro subtitulo el prompt que se utilizo. Todo este manual generar en formato latex"}
\end{quote}

\end{document}












